% Slides with algoviz figures. Compile with shell escape:
%   latexmk -pdf -shell-escape slides.tex
% steps= makes one slide per frame of a series, all in one picture.
% Inside a frame, write the # of a step as \#, or make the frame fragile.
% A figure fills the line and three quarters of the frame's height; height=
% gives another share.
\documentclass{beamer}
\usepackage{algoviz}

\begin{document}

% steps without a value is every frame of the run.
\begin{frame}{Dijkstra's algorithm, step by step}
  \centering
  \algovizfile[algorithm=dijkstra, steps]{graph.json}
\end{frame}

% A range of steps. After \pause, the series starts on the next slide.
\begin{frame}{Settling $a$}
  The queue hands out $a$, at distance $3$, and relaxes its edges.
  \pause

  \centering
  \algovizfile[algorithm=dijkstra, steps=extract\#3..improve\#5,
    height=0.6\textheight]{graph.json}
\end{frame}

% step= is one figure, on every slide of its frame.
\begin{frame}[fragile]{Dijkstra's algorithm: the last relaxation}
  \centering
  \algovizfile[algorithm=dijkstra, step=relax#6]{graph.json}
\end{frame}

% A section of a run, in the print theme.
\begin{frame}{Deleting $3$ from a binary search tree}
  \centering
  \algoviz[steps=delete-3, theme=print]{bst.delete}{
    "from": [8, 3, 10, 1, 6, 4, 7], "key": 3}
\end{frame}

\end{document}
